% Capitolo 1: Stato dell'arte
\chapter{Stato dell'arte}
\label{chap:stato_arte}

\section{Introduzione alle applicazioni mobili}
Le applicazioni mobili sono software progettati per dispositivi portatili (smartphone, tablet) e si caratterizzano per vincoli di risorse (batteria, CPU, memoria), modalità di interazione basate sul touch e requisiti variabili di connettività. Negli ultimi anni gli smartphone sono diventati il principale canale di fruizione di servizi personali e di gruppo, con crescenti esigenze in termini di usabilità, sincronizzazione dei dati e funzionamento in assenza di rete (offline).

\section{Paradigmi e categorie}
Le principali categorie tecnologiche e di sviluppo per applicazioni mobili includono:

\begin{itemize}
  \item \textbf{Native}: applicazioni sviluppate con SDK e linguaggi nativi (Swift per iOS, Kotlin/Java per Android). Offrono massimo controllo e prestazioni, ma richiedono sviluppi distinti per ciascuna piattaforma.
  \item \textbf{Cross‑platform}: framework quali React Native, Flutter e soluzioni basate su Expo consentono di condividere ampie porzioni di codice tra i sistemi, riducendo tempi e costi di sviluppo a fronte di possibili limitazioni nell'accesso ad API native avanzate.
  \item \textbf{Progressive Web App (PWA)}: applicazioni web con comportamento "app‑like" installabili da browser. Offrono distribuzione semplice ma limitazioni nell'accesso ad alcune funzionalità hardware e in alcuni casi prestazioni inferiori rispetto a soluzioni native.
\end{itemize}

\section{Tecnologie e componenti rilevanti}
Per la tipologia di applicazione trattata in questa tesi, le tecnologie rilevanti comprendono:
\begin{itemize}
  \item \textbf{React Native / Expo}: rapido ciclo di sviluppo in JavaScript/TypeScript e ampio ecosistema di librerie.
  \item \textbf{Flutter}: UI declarative e buone prestazioni su entrambe le piattaforme mobili.
  \item \textbf{DB on‑device}: SQLite (tramite wrapper), Realm — scelte fondamentali per resilienza offline.
  \item \textbf{BaaS e servizi cloud}: Firebase, Supabase per autenticazione, realtime e storage.
  \item \textbf{Strumenti di testing e analytics}: servizi per raccolta crash, metriche d'uso e test di usabilità remota o in laboratorio.
\end{itemize}

\section{Aree applicative affini}
Nel dominio della gestione sportiva e della amministrazione di piccoli gruppi, le funzionalità ricorrenti sono: gestione anagrafica di team e giocatori, registrazione delle presenze, annotazione di sanzioni ("multe") e penalità, canali di comunicazione (chat/annunci) e funzionalità di reporting. Diverse soluzioni commerciali e open source coprono alcuni di questi aspetti, ma raramente offrono una combinazione completa con forte attenzione all'affidabilit\'a offline e a flussi semplificati per utenti non tecnici.

\section{Esempi di soluzioni esistenti}
Di seguito si fornisce una breve panoramica qualitativa di categorie di prodotti che possono considerarsi competitor o affini:

\begin{itemize}
  \item \textbf{Applicazioni per tracking presenze}: possono offrire funzionalit\'a solide per le presenze ma tradizionalmente dipendono dalla connettivit\'a e offrono limitata persistenza locale.
  \item \textbf{App di comunicazione e calendario per team}: forniscono ottime funzionalit\'a di messaggistica e pianificazione, ma spesso non includono workflow dedicati alla gestione di multe/penalit\'a o reporting esportabile.
  \item \textbf{Soluzioni ERP o gestionali web}: ricche di funzionalit\'a ma spesso complesse e non ottimizzate per interazioni mobili rapide.
\end{itemize}

\section{Analisi comparativa}
Si suggerisce di includere una tabella comparativa che metta a confronto le caratteristiche chiave delle soluzioni esistenti e del lavoro proposto. Qui sotto un esempio di tabella da adattare ai casi concreti e ai nomi delle app analizzate.

\begin{table}[H]
\centering
\begin{tabular}{|l|c|c|c|c|}
\hline
\textbf{Caratteristica} & \textbf{App A} & \textbf{App B} & \textbf{App C} & \textbf{Questo lavoro} \\ \hline
Autenticazione & \checkmark & \checkmark & \checkmark & \checkmark (ruoli) \\ \hline
Persistenza offline & -- & parziale & -- & \checkmark (SQLite) \\ \hline
Sincronizzazione cloud & parziale & \checkmark & parziale & \checkmark (resiliente) \\ \hline
Chat interna & -- & \checkmark & \checkmark & \checkmark (essenziale) \\ \hline
Reporting/esportazione & parziale & -- & \checkmark & \checkmark (CSV/XLSX) \\ \hline
Facilit\'a d'uso & media & alta & bassa & alta (role‑based) \\ \hline
\end{tabular}
\caption{Schema comparativo esemplificativo}
\label{tab:comparativa}
\end{table}

\section{Limiti e criticit\'a delle soluzioni esistenti}
Dall'analisi della letteratura tecnica e delle app emerse si riscontrano criticit\'a ricorrenti:
\begin{itemize}
  \item Dipendenza dalla connettivit\'a: molte app perdono funzionalit\'a essenziali in assenza di rete.
  \item Strategie di sincronizzazione deboli: assenza di meccanismi affidabili per risolvere conflitti e per retry in caso di errori di rete.
  \item UX complessa per utenti non tecnici: flussi non ottimizzati per operazioni frequenti (es. applicare una multa).
  \item Mancanza di un approccio role‑based semplice: permessi e interfacce non sempre differenziati per ruoli come "mister" e "player".
\end{itemize}

\section{Motivazioni e contributi del lavoro}
Questo lavoro si pone l'obiettivo di colmare i gap sopra elencati offrendo i seguenti contributi:
\begin{enumerate}
  \item \textbf{Approccio offline‑first}: persistenza locale tramite SQLite per garantire disponibilit\'a delle funzionalit\'a critiche anche senza rete.
  \item \textbf{Sincronizzazione resiliente}: integrazione con un backend leggero (es. Supabase) per sincronizzazione incrementale, gestione dei conflitti e meccanismi di retry.
  \item \textbf{UX role‑based e semplificata}: interfacce distinte e ottimizzate per i ruoli principali (mister, player) per semplificare le operazioni pi\`u frequenti.
  \item \textbf{Architettura modulare e testabile}: separazione tra UI, logica e livello DB mediante hook e layer DB per facilitare manutenzione e test.
  \item \textbf{Validazione tramite user testing}: progettazione guidata da user stories e valutazione empirica dei flussi tramite test utente con metriche quantitative e qualitative.
\end{enumerate}

\section{Riferimenti e letteratura consigliata}
Per completare la sezione bibliografica si consigliano riferimenti su temi quali: sincronizzazione dati in ambiente mobile (offline‑first, CRDTs), studi di usabilità su applicazioni sportive/gestionali, confronti fra framework cross‑platform e native e documentazione tecnica degli strumenti utilizzati (React Native/Expo, Supabase, SQLite). Inserire riferimenti formali nella bibliografia (file \texttt{refs.bib}).

\section{Figure e tabelle da includere}
Consiglio di aggiungere le seguenti figure/tabelle nel capitolo o nei capitoli successivi:
\begin{itemize}
  \item Tabella comparativa (\autoref{tab:comparativa}).
  \item Diagramma architetturale ad alto livello (UI $\leftrightarrow$ hooks $\leftrightarrow$ DB locale $\leftrightarrow$ backend).
  \item Flussi utente per i ruoli principali (mister, player).
  \item Screenshot o wireframe dell'app (inseriti in Cap. 2/3 quando disponibili).
\end{itemize}

\section{Conclusione}
In sintesi, pur essendo disponibili molte soluzioni che affrontano singoli aspetti della gestione di team e comunicazione, rimane una lacuna nell'integrazione di un'esperienza mobile semplice e affidabile che combini resilienza offline, sincronizzazione robusta e un'interfaccia role‑based accessibile. Questa tesi propone un'applicazione orientata a colmare tale gap e a validare empiricamente l'efficacia delle scelte progettuali.

% Fine Capitolo 1
